\documentclass{situationist-preprint}

\title{Partial Unity in Split-Brain Cases}
\msid{MND-15-a}
\msfamily{MND-15}
\msdiscipline{Philosophy of Mind}
\mstype{Principal manuscript}
\mskeywords{unity of consciousness; split-brain; commissurotomy; co-consciousness; partial unity; streams of consciousness}

\begin{document}
\maketitle

\begin{abstract}
Is consciousness unified? In normal waking subjects, all experiences had at a time seem to belong to a single total experience. Split-brain patients, whose cerebral hemispheres have been surgically disconnected, appear under laboratory conditions to have experiences that are not integrated with one another, and the standard interpretations say that they have two streams of consciousness, or one stream that switches between hemispheres. We argue that consciousness is unified at each time in normal subjects, and that split-brain cases involve partial unity rather than two streams. Unity is a relation of co-consciousness between experiences, and partial unity is the case in which this relation fails to be transitive: two experiences are each co-conscious with a third without being co-conscious with each other. We show that the evidence from split-brain patients, including the recent finding that they can respond to stimuli in either visual field with either hand, is better explained by partial unity than by the two-streams or the switch model, and we answer the objection that partial unity is inconceivable. The model extends to the distributed nervous system of the octopus and fixes the sense in which partial hallucination is unified.
\end{abstract}

\section{The problem}

When I hear a car pass while seeing the page in front of me and feeling the chair beneath me, I do not have three experiences that happen to be simultaneous; I have one experience that includes all three. The \emph{unity thesis} says that this holds generally: at each time, all the experiences of a normal subject are parts of a single total experience \citep{BayneChalmers2003,Bayne2010}. It has been held as a datum by philosophers since Descartes and Kant and is presupposed by the ordinary practice of counting subjects by counting bodies. MND-06 uses it as a criterion for what counts as a single mind.

Patients whose corpus callosum has been severed to relieve epilepsy pose the main challenge. In everyday life they act as unified agents. But in experiments in which a stimulus is flashed to one visual hemifield only, so that it reaches only the opposite hemisphere, they can name a word presented to the right hemifield, processed by the language-dominant left hemisphere, but not one presented to the left; and they can pick out with the left hand, controlled by the right hemisphere, an object matching a picture presented to the left hemifield, while denying verbally that they saw anything \citep{Sperry1968,Gazzaniga2005}. \citet{Nagel1971} argued that none of the available answers to \enquote{how many minds?} is acceptable: not one, since the hemispheres' experiences are not integrated; not two, since in everyday life the patient acts as one; and not a number that changes with the experimental setting. He concluded that the concept of a single subject may not apply cleanly.

Four interpretations have been proposed. On the \emph{two-streams} model, the patient has two streams of consciousness at all times, or at least under experimental conditions, one associated with each hemisphere \citep{Sperry1968}; \citet{Schechter2018} has developed a sophisticated version on which there are two subjects of experience who share a single body and, in an extended sense, a single self. On the \emph{switch} model, the patient has a single stream at each time, but its contents switch between the hemispheres according to which is engaged by the task \citep{Bayne2008}. On the \emph{partial unity} model, there is one subject whose experiences are not all co-conscious, but whose co-consciousness relation fails transitivity \citep{Lockwood1989}. Fourth, Nagel's no-answer view. We defend the partial unity model, against the objection, pressed by \citet{Hurley1998} and \citet{Bayne2008}, that it is inconceivable.

\section{Framework}

\begin{definition}[Co-consciousness]\label{d:co}
Two experiences $e_1$ and $e_2$ had at the same time are \emph{co-conscious} if there is an experience of which both are parts, that is, if there is something it is like to have them together over and above what it is like to have each \citep{BayneChalmers2003,Dainton2000}.
\end{definition}

\begin{definition}[Unity; partial unity; streams]\label{d:unity}
A subject's experiences at a time are \emph{unified} if every two of them are co-conscious. They are \emph{partially unified} if the co-conscious relation among them is connected, in that any two are linked by a chain of co-conscious pairs, but not every two are co-conscious. A \emph{stream} at a time is a maximal set of experiences linked by chains of co-conscious pairs.
\end{definition}

On Definition~\ref{d:unity}, unity and partial unity both yield one stream; two streams require two sets of experiences with no co-conscious chain between them. Co-consciousness is reflexive and symmetric; partial unity is the case in which it is not transitive. The companion manuscript MND-15-b develops the method of counting streams that these definitions induce.

\begin{principle}[Integration]\label{p:int}
If the contents of two simultaneous experiences of a subject are jointly available to her for the flexible control of a single action, comparison, or report, the experiences are co-conscious; if two experiences share a common experienced part, such as a bodily sensation or an emotion, they are co-conscious.
\end{principle}

The first clause is the evidential principle used by every party to the debate: integration of content in the control of action and report is how co-consciousness is diagnosed, in others and in the normal case. The second follows from Definition~\ref{d:co}: if one token experience is a part of both $e_1$ and $e_2$, then $e_1$ and $e_2$ overlap in a single experience of which each is a part together with it.

\begin{principle}[Parsimony of tokens]\label{p:pars}
Experience tokens are not to be multiplied beyond what the evidence of integration and dissociation requires.
\end{principle}

\section{Unity in normal subjects}

\begin{lemma}\label{l:normal}
In normal waking subjects, consciousness is unified at each time.
\end{lemma}

\begin{proof}
In normal subjects, the contents of the experiences of different modalities and of different parts of the visual field are jointly available for comparison, report and the control of action: one can say that the sound came from the left of what one sees, compare the color at two points of the field, reach for a seen object while hearing an instruction. By Principle~\ref{p:int}, every two experiences whose contents are so available are co-conscious. No dissociation in normal subjects shows two simultaneous experiences whose contents are not jointly available when the subject attends to both. So every two of a normal subject's experiences at a time are co-conscious, and by Definition~\ref{d:unity} they are unified.
\end{proof}

The restriction to experiences whose contents are available when the subject attends to both is not a restriction to attended experiences: MND-24 argues that consciousness extends beyond attention, and the unity claim concerns all experiences had at a time, since any of them could be brought into comparison with any other.

\section{The split-brain evidence}

Three bodies of evidence matter. First, under lateralized presentation, visual contents in the two hemifields are not jointly available for comparison or report: the patient cannot say whether two stimuli flashed to the two hemifields are the same \citep{Gazzaniga2005}. Second, many contents are shared. The patients have a single body image and bodily sensations that are available to both hemispheres through ipsilateral and subcortical pathways; emotional reactions produced by a stimulus presented to one hemisphere are felt and reported by the other; auditory input reaches both hemispheres. Third, \citet{Pinto2017} found that two patients could indicate the presence and location of stimuli in either visual field with either hand and verbally, although they could not compare stimuli across the fields. Perception was divided; awareness of presence and location was not.

\begin{lemma}\label{l:not2}
The evidence does not support two streams.
\end{lemma}

\begin{proof}
Let $L$ and $R$ be the patient's visual experiences of the left and right hemifields under lateralized presentation, and $B$ her bodily experience at the same time. By the first body of evidence and Principle~\ref{p:int}, there is no evidence that $L$ and $R$ are co-conscious, and we grant that they are not. By the second, the contents of $B$ are jointly available with those of $L$ in the control of the left hand, and with those of $R$ in verbal report; so $B$ is co-conscious with each of $L$ and $R$. Two streams would require that no chain of co-consciousness links $L$ and $R$, so that $B$ would have to be duplicated into two token bodily experiences, one in each stream, with the same content. Nothing in the evidence calls for two tokens rather than one: there is a single body image, and the duplicated experiences would never diverge in content. By Principle~\ref{p:pars}, the duplication is not warranted. The third body of evidence adds that awareness of presence and location in each hemifield is available to both response systems, which on the two-streams model would require a further duplication.
\end{proof}

\begin{lemma}\label{l:noswitch}
The evidence does not support the switch model.
\end{lemma}

\begin{proof}
On the switch model, at any time only one hemisphere's visual contents are conscious. But in bimanual tasks, split-brain patients respond simultaneously and appropriately with each hand to stimuli in the corresponding hemifield, and in the experiments of \citet{Pinto2017} each response system had access to stimuli in both fields. If integration in the control of action is the criterion of consciousness, as all parties accept, then both hemifield contents are conscious at the same time, which the switch model denies.
\end{proof}

\begin{theorem}\label{t:main}
Consciousness is unified at each time in normal subjects. In split-brain patients, it is partially unified: there is a single stream, in which some pairs of experiences are not co-conscious.
\end{theorem}

\begin{proof}
The first sentence is Lemma~\ref{l:normal}. For the second, take $L$, $R$ and $B$ as in Lemma~\ref{l:not2}. $L$ and $R$ are not co-conscious; each is co-conscious with $B$. By Lemma~\ref{l:not2}, $B$ is a single token. So the co-conscious relation among the patient's experiences is connected but not transitive. By Definition~\ref{d:unity}, the experiences are partially unified and form one stream. By Lemma~\ref{l:noswitch}, they are simultaneous, so the switch model does not describe them. Nagel's no-answer view is not needed, since the question how many streams there are has an answer: one, partially unified.
\end{proof}

\section{Objections and replies}

\begin{objection}[Inconceivability]\label{o:inconc}
We cannot imagine what it would be like to have two experiences each co-conscious with a third but not with each other. Since what it is like to be in a total state must be imaginable from the subject's point of view, partial unity is impossible \citep{Hurley1998,Bayne2008}.
\end{objection}

\begin{reply}
To imagine a subject's experience from her point of view is to imagine a single total experience, and partial unity is precisely the hypothesis that there is no single total experience including all of her experiences. So partial unity predicts that it cannot be imagined as a whole, and its unimaginability is no evidence against it. What can be imagined is each maximal unified part: what it is like to have $L$ with $B$, and what it is like to have $R$ with $B$. Nothing in the concept of co-consciousness requires transitivity; it is a relation of joint inclusion in an experience, and two experiences each jointly included with a third need not be jointly included with each other, any more than two overlapping regions need contain each other's parts \citep{Lockwood1989}.
\end{reply}

\begin{objection}[Subject]\label{o:subject}
A single subject is a single point of view, and a point of view is unified by definition. Partial unity therefore implies two subjects.
\end{objection}

\begin{reply}
The objection defines the subject by unity and so begs the question. The patient is one subject in every respect that the evidence bears on: she has one body image, one emotional life, one set of memories and plans, and acts as one agent in ordinary life. What she lacks, under special conditions, is co-consciousness between some visual contents. The two-subjects view of \citet{Schechter2018} reconciles this with two subjects by distinguishing subjects of experience from persons, but at the cost of duplicating the shared experiences, which Lemma~\ref{l:not2} rejects.
\end{reply}

\begin{objection}[Access]\label{o:access}
Principle~\ref{p:int} diagnoses access, not phenomenal co-consciousness. $L$ and $R$ might be phenomenally co-conscious although their contents are not jointly accessible.
\end{objection}

\begin{reply}
If so, the patient's consciousness is fully unified, and the conclusion that there are not two streams holds a fortiori. The principle is used only to deny co-consciousness where integration is absent; if the denial is too strong, the result moves toward unity, not away from it.
\end{reply}

\begin{objection}[Ad hoc]\label{o:adhoc}
Partial unity is an artefact of the experiments; outside them the patient is unified, so the model applies only intermittently.
\end{objection}

\begin{reply}
That is the model's prediction, not a difficulty. Co-consciousness between $L$ and $R$ fails only when their contents are confined to different hemispheres, which in daily life eye movements and shared input prevent. The model describes one stream at all times, more or less unified according to the input, as the evidence requires.
\end{reply}

\section{Further applications}

\paragraph{Octopuses.} MND-23 concludes that cephalopods are probably conscious, with distributed processing: about two thirds of an octopus's neurons are in its arms, which perform complex movements under local control. Theorem~\ref{t:main} supplies a model. The experiences associated with different arms may fail to be co-conscious with one another while each is co-conscious with the central experience of the animal's body and surroundings that the central brain sustains. The octopus would then have a single, partially unified stream, rather than nine subjects or none, with more partiality than a split-brain patient under laboratory conditions. Whether this is so is a question about integration that Principle~\ref{p:int} makes empirical.

\paragraph{Partial hallucination.} MND-16 holds that hallucinations are experiences of a different kind from veridical perceptions, and that in partial hallucination a perception of a real table and an imagistic experience of a cup are integrated into one visual field. The framework of this manuscript fixes what that integration is: the two components are co-conscious in the sense of Definition~\ref{d:co}, and co-consciousness is a relation between experiences that does not require them to be of a common kind. In a normal subject, partial hallucination is therefore a case of full unity across components of different kinds; were the integration to fail, as it may in some pathological cases in which the subject cannot relate the hallucinated object to the perceived scene, the case would be one of partial unity on the split-brain model.

\section{What the result does not establish}

The argument does not establish that consciousness is unified across time, nor that the unity thesis holds in states such as dreaming, anesthesia or dissociative disorders, which raise separate questions. It does not establish that co-consciousness is identical with any neural relation. And it rests on the evidence from a small number of patients, whose surgery and history vary; if a patient were found in whom no experience is shared between the hemispheres, the theorem would classify that patient as having two streams.

\section{Conclusion}

Normal consciousness is unified: every two experiences a normal subject has at a time are parts of a single experience. Split-brain patients are not counterexamples to the existence of a single stream. Their visual experiences in the two hemifields may fail to be co-conscious, but each is co-conscious with bodily, emotional and other shared experiences, and the patient has one stream, partially unified. The model is conceivable in the only sense that matters, and it extends to other minds whose processing is distributed.

\begin{thebibliography}{99}\small
\bibitem[Bayne(2008)]{Bayne2008} Bayne, T. (2008). The unity of consciousness and the split-brain syndrome. \emph{The Journal of Philosophy}, 105(6), 277--300.
\bibitem[Bayne(2010)]{Bayne2010} Bayne, T. (2010). \emph{The Unity of Consciousness}. Oxford University Press.
\bibitem[Bayne and Chalmers(2003)]{BayneChalmers2003} Bayne, T. and Chalmers, D.~J. (2003). What is the unity of consciousness? In A. Cleeremans (ed.), \emph{The Unity of Consciousness: Binding, Integration, and Dissociation}. Oxford University Press.
\bibitem[Dainton(2000)]{Dainton2000} Dainton, B. (2000). \emph{Stream of Consciousness: Unity and Continuity in Conscious Experience}. Routledge.
\bibitem[Gazzaniga(2005)]{Gazzaniga2005} Gazzaniga, M.~S. (2005). Forty-five years of split-brain research and still going strong. \emph{Nature Reviews Neuroscience}, 6(8), 653--659.
\bibitem[Hurley(1998)]{Hurley1998} Hurley, S.~L. (1998). \emph{Consciousness in Action}. Harvard University Press.
\bibitem[Lockwood(1989)]{Lockwood1989} Lockwood, M. (1989). \emph{Mind, Brain and the Quantum: The Compound \enquote{I}}. Blackwell.
\bibitem[Nagel(1971)]{Nagel1971} Nagel, T. (1971). Brain bisection and the unity of consciousness. \emph{Synthese}, 22(3--4), 396--413.
\bibitem[Pinto et~al.(2017)]{Pinto2017} Pinto, Y., Neville, D.~A., Otten, M., Corballis, P.~M., Lamme, V.~A.~F., de~Haan, E.~H.~F., Foschi, N. and Fabri, M. (2017). Split brain: Divided perception but undivided consciousness. \emph{Brain}, 140(5), 1231--1237.
\bibitem[Schechter(2018)]{Schechter2018} Schechter, E. (2018). \emph{Self-Consciousness and \enquote{Split} Brains: The Minds' I}. Oxford University Press.
\bibitem[Sperry(1968)]{Sperry1968} Sperry, R.~W. (1968). Hemisphere deconnection and unity in conscious awareness. \emph{American Psychologist}, 23(10), 723--733.
\end{thebibliography}
\end{document}
